\section{Reproducibility Findings}
\label{sec:reproducibility}

% Target: 2 pages (~1600 words). Heaviest evidence section. Source:
% GENERATION_FAILURE_ANALYSIS.md (now with file:line citations) and
% paper_underspecifications.md.

\subsection{Three implementation choices in the released code}
\label{sec:repro:choices}

% Each subsubsection cites file:line. Pull verbatim from
% GENERATION_FAILURE_ANALYSIS.md §Provenance table.

\paragraph{Deterministic argmax in the PEM rollout.}
% The released sampling code rolls out the PEM with torch.argmax for pattern
% and length heads, and the raw regression mean for magnitude — no sampling.
% Cite: code_from_authors/codes/models/sampling.py lines 30, 33, 34. Quote
% the three lines verbatim in a listing. Show that an autoregressive rollout
% under a deterministic map must converge to a fixed point or short cycle.

\paragraph{Decoder hidden dimension of one.}
% pgm_params['sae_hidden_dim'] = 1 in the released config; the Scaling-AE
% decoder is nn.LSTM(hidden_dim, hidden_dim, 2). One-unit LSTM hidden state.
% Cite: code_from_authors/codes/models/model_params.py line 33;
% scaling_autoencoder.py line 58.

\paragraph{Magnitude range-normalization commented out.}
% The sampling-time generation step has the range-normalization that would
% enforce |segment| = β commented out. Cite: sampling.py lines 53–55.

\todo{For each of the three choices, show one short listing of the relevant
code lines (use the \texttt{listings} package). Anchor each to a file:line
reference. Keep listings under 5 lines each.}

\subsection{Observed degeneracy}
\label{sec:repro:degeneracy}

% Quantitative + visual evidence:
% - All 30 saved runs per (asset, k, protocol) converge to near-identical
%   trajectories; show std-across-runs ≈ 0.
% - PEM state-trajectory analysis: S&P 500 collapses to fixed point (pattern
%   12 used 99.6%); GOOG and ZC=F collapse to period-2 cycles.
% - PGM decoder output is near-constant across all 14 patterns (cross-pattern
%   std < 0.001).

\begin{figure}[t]
  \centering
  \todo{Figure: composite of (a) S\&P\,500 synthetic trajectory (linear ramp);
  (b) PEM state sequence (pattern $\in\{12\}$ for 2519/2520 steps);
  (c) PGM decoder output across 14 patterns showing near-identical curves.}
  \caption{Diagnostic plots of the three failure modes. \todo{Replace.}}
  \label{fig:degeneracy}
\end{figure}

\begin{table}[t]
  \centering
  \todo{Table: per-asset summary of synthetic trajectory degeneracy. Columns:
  asset, K, dominant pattern share, std across 30 runs at day 252, std across
  30 runs at day 2500. Numbers from \texttt{scripts/investigate\_pgm.py} and
  the saved \texttt{synthetic/*/run\_*\_states.npy}.}
  \caption{Degeneracy summary across assets. \todo{Fill in numbers from E1
  pipeline.}}
  \label{tab:degeneracy_summary}
\end{table}

\subsection{H1: full-budget retraining}
\label{sec:repro:h1}

% Per GENERATION_FAILURE_ANALYSIS.md §Alternative hypotheses / H1.
% The released config sets PGM n_epochs=400 and PEM n_epochs=1000 (we used 30
% and 60 for the artifacts above). To rule out undertraining as the cause,
% we retrain at the authors' full epoch budget on S&P 500.
%
% Outcome (E5):
% \todo{Fill in after E4/E5 complete. State precisely whether the decoder
% cross-pattern std exceeds the 0.01 falsification threshold. If yes:
% undertraining was the cause; revise the framing of §4 accordingly. If no:
% the degeneracy is robust to epoch budget, strengthening the reproducibility
% claim.}

\subsection{Underspecifications between paper and code}
\label{sec:repro:underspecs}

% Source: paper_underspecifications.md
% Pull the 11 entries into a single compact table:
% column 1: design choice; column 2: paper specifies?; column 3: released
% value; column 4: notes (e.g., "potential confound").
% Move the full text version to the appendix.

\begin{table}[t]
  \centering
  \todo{Table: underspecifications. Rows: each of the 11 entries from
  \texttt{paper\_underspecifications.md}. Keep this table at 6-8 rows in the
  main text; full table in appendix.}
  \caption{Paper-vs-code underspecifications and our choices.}
  \label{tab:underspecs}
\end{table}
